```latex
\documentclass[11pt]{article}
\usepackage{amsmath,amssymb,amsthm}
\usepackage[colorlinks=true,urlcolor=blue]{hyperref}

\newtheorem{theorem}{Theorem}
\newtheorem{lemma}{Lemma}
\newtheorem{definition}{Definition}
\theoremstyle{remark}
\newtheorem{remark}{Remark}

\title{Local $J$-holomorphic equations for symplectic simple normal crossings divisors}
\author{}
\date{}

\begin{document}
\maketitle

\begin{abstract}
We give a complete proof that, after passing to a cohomologous symplectic form along a path preserving a given symplectic simple normal crossings divisor, there exists a tame almost complex structure for which the divisor admits local smooth complex coordinates with each defining function $J$-holomorphic.
\end{abstract}

\section{Statement of the result}

Let $(M,\omega)$ be a closed symplectic manifold and let
\[
D=\bigcup_{i=1}^{N}D_i\subset M
\]
be a symplectic simple normal crossings divisor. We take the standard definition: each $D_i$ is a closed symplectic submanifold of real codimension $2$, all intersections
\[
D_I:=\bigcap_{i\in I}D_i,\qquad I\subset \{1,\dots,N\},
\]
are transverse and symplectic, and near every point of $D_I$ there are Darboux coordinates in which the branches $D_i$ are coordinate hyperplanes. See \cite{FZ17} for this definition and for the associated space of compatible almost complex structures.

The question is whether one can find an $\omega$-tame almost complex structure $J$ such that, locally on $M$, there are smooth complex coordinates $(z_1,\dots,z_n)$ with
\[
D=\{z_1\cdots z_r=0\}
\]
for some $0\le r\le n$, and such that $z_1,\dots,z_r$ are $J$-holomorphic. As the question allows, we may first replace $\omega$ by a symplectic form joined to it by a path of cohomologous symplectic forms along which $D$ remains a symplectic simple normal crossings divisor.

The answer is affirmative.

\begin{theorem}\label{thm:main}
Let $(M,\omega)$ be a closed symplectic manifold and $D\subset M$ a symplectic simple normal crossings divisor. Then there exists a smooth path
\[
\{\omega_t\}_{t\in[0,1]},\qquad [\omega_t]=[\omega]\in H^2_{\mathrm{dR}}(M),
\]
such that $D$ is a symplectic simple normal crossings divisor for every $\omega_t$, and an $\omega_1$-compatible almost complex structure $J$ on $M$ with the following property. For every $p\in M$, letting
\[
I(p):=\{i\in\{1,\dots,N\}:p\in D_i\},\qquad r:=|I(p)|,
\]
there are smooth complex coordinates $(z_1,\dots,z_n)$ on a neighborhood of $p$ such that
\[
D\cap U=\{z_1\cdots z_r=0\},
\]
and the functions $z_1,\dots,z_r$ are $J$-holomorphic:
\[
dz_i(Jv)=i\,dz_i(v)
\]
for every tangent vector $v$.
\end{theorem}

Thus, in the sense allowed by the question, such a $J$ exists. Since an $\omega_1$-compatible almost complex structure is in particular $\omega_1$-tame, the theorem answers the question affirmatively.

\section{Regularizations and almost K\"ahler data}

The proof uses the regularization theorem for symplectic simple normal crossings divisors due to Farajzadeh Tehrani--McLean--Zinger. We recall the part of that theory needed here; the full definitions are in \cite{TMZ21}, and the existence theorem is proved in \cite{TMZ18} and \cite{NSF}.

\begin{definition}[Almost K\"ahler data]\label{def:AK}
Let $D\subset(M,\omega)$ be a symplectic simple normal crossings divisor. An \emph{$\omega$-regularization} $\mathcal R$ of $D$ consists, in particular, of the following data. For each $I\subset\{1,\dots,N\}$, let
\[
\mathcal N_I:=\bigoplus_{i\in I}N_MD_i\big|_{D_I}
\]
be the symplectic normal bundle of $D_I$. The regularization gives a complex structure on each normal line bundle $N_MD_i$, a neighborhood $\mathcal O_I\subset\mathcal N_I$ of the zero section, and a diffeomorphism
\[
\Psi_I:\mathcal O_I\to U_I\subset M
\]
onto a tubular neighborhood of $D_I$. The splitting
\[
T\mathcal O_I=\pi_I^*TD_I\oplus \bigoplus_{i\in I}N_MD_i\big|_{D_I}
\]
is $\omega$-orthogonal, where $\pi_I:\mathcal N_I\to D_I$ is the bundle projection.

An almost complex structure $J$ is \emph{$\mathcal R$-compatible} if the triple $(\omega,\mathcal R,J)$ belongs to the space $AK(M,D)$ of almost K\"ahler data defined in \cite{TMZ18}. In particular, for every $I$,
\begin{enumerate}
\item $D_I$ is $J$-holomorphic;
\item on each fiber $F$ of $\pi_I\big|_{\mathcal O_I}$, the pullback of $J$ by $\Psi_I$ is the standard fiberwise complex structure on $\bigoplus_{i\in I}\mathbb C$;
\item the map
\[
\pi:=\pi_I\circ \Psi_I^{-1}:U_I\to D_I
\]
is $(J,J|_{D_I})$-holomorphic, i.e.
\[
d\pi\circ J=J|_{D_I}\circ d\pi.
\]
\end{enumerate}
\end{definition}

The existence of such data is the main point of the regularization theorem.

\begin{theorem}[Regularization existence, {\cite[Theorem 7]{TMZ18}}]\label{thm:reg}
Let $(M,\omega)$ be a closed symplectic manifold and $D\subset M$ a symplectic simple normal crossings divisor. Then there exists a smooth path $\{\omega_t\}_{t\in[0,1]}$ of cohomologous symplectic forms on $M$, with $\omega_0=\omega$, such that $D$ is a symplectic simple normal crossings divisor with respect to each $\omega_t$, and such that $D$ admits an $\omega_1$-regularization $\mathcal R$. Moreover, there exists an $\omega_1$-compatible, $\mathcal R$-compatible almost complex structure $J$ on $M$; in the notation of \cite{TMZ18},
\[
(\omega_1,\mathcal R,J)\in AK(M,D).
\]
\end{theorem}

\begin{remark}
The theorem is proved in \cite{TMZ18} and is also presented in \cite{NSF}. The construction of the space $AK(M,D)$ and the proof that the forgetful projection from $AK(M,D)$ to the auxiliary data is a weak homotopy equivalence imply the existence of the compatible $J$ above.
\end{remark}

\section{Local $J$-holomorphic defining functions}

We now prove the key local lemma: an $\mathcal R$-compatible $J$ gives rise to the desired coordinates.

\begin{lemma}\label{lem:local}
Let $(\omega,\mathcal R,J)\in AK(M,D)$. For every $p\in D$, with
\[
I:=\{i:p\in D_i\},\qquad r:=|I|,
\]
there exist a neighborhood $U$ of $p$ and smooth complex coordinates
\[
(z_1,\dots,z_r,w_1,\dots,w_{n-r})
\]
on $U$ such that
\[
D_i\cap U=\{z_i=0\}\quad (1\le i\le r),
\]
and each $z_i$ is $J$-holomorphic. In particular,
\[
D\cap U=\{z_1\cdots z_r=0\}.
\]
\end{lemma}

\begin{proof}
Let $p\in D_I$. By Definition~\ref{def:AK}, the $\omega$-regularization gives a neighborhood $\mathcal O_I$ of the zero section in
\[
\mathcal N_I=\bigoplus_{i\in I}N_MD_i\big|_{D_I},
\]
together with a diffeomorphism
\[
\Psi_I:\mathcal O_I\to U_I\subset M.
\]
Shrinking $U_I$ if necessary, we may assume that $U_I$ is disjoint from every $D_j$ with $j\notin I$.

Choose a neighborhood $W$ of $p$ in $D_I$ and smooth complex coordinates $(w_1,\dots,w_{n-r})$ on $W$. Since each $N_MD_i$ is a complex line bundle, shrinking $W$ if necessary we may trivialize
\[
N_MD_i\big|_W\cong W\times\mathbb C
\]
and denote the fiber coordinate by $z_i$. Thus, on
\[
\mathcal O_I\cap \pi_I^{-1}(W),
\]
we have smooth complex coordinates
\[
(z_1,\dots,z_r,w_1,\dots,w_{n-r}).
\]
Because $\Psi_I$ is a diffeomorphism, these pull back to smooth complex coordinates on
\[
U:=\Psi_I\bigl(\mathcal O_I\cap\pi_I^{-1}(W)\bigr)\subset U_I.
\]
We abuse notation and denote the pulled-back functions again by $z_i,w_j$.

By construction,
\[
D_i\cap U=\{z_i=0\}\qquad (1\le i\le r),
\]
and therefore
\[
D\cap U=\bigcup_{i=1}^r D_i\cap U=\{z_1\cdots z_r=0\}.
\]

It remains to show that each $z_i$ is $J$-holomorphic. Let
\[
\pi:=\pi_I\circ\Psi_I^{-1}:U\to D_I
\]
and let
\[
V:=\ker d\pi\subset TU
\]
be the vertical tangent subbundle. In the coordinates above, $V$ is the tangent bundle of the fibers of the normal bundle. By Definition~\ref{def:AK}(2), $J$ preserves $V$ and, on each fiber, is the standard complex structure on $\mathbb C^r$. Hence, for $v\in V$,
\[
dz_i(Jv)=i\,dz_i(v).
\]

Since $\mathcal R$ is an $\omega$-regularization, the horizontal subbundle
\[
H:=V^{\perp_\omega}
\]
is precisely the tangent bundle to the base $D_I$ in the product decomposition. In particular,
\[
dz_i\big|_H=0.
\]
Moreover, because $J$ is $\omega$-compatible and $J(V)=V$, we have $J(H)=H$. Indeed, if $h\in H$ and $v\in V$, then
\[
\omega(Jh,v)=\omega(Jh,J(-Jv))=\omega(h,-Jv)=0,
\]
so $Jh\in V^{\perp_\omega}=H$.

Now every tangent vector $v\in TU$ decomposes as
\[
v=v_H+v_V,\qquad v_H\in H,\quad v_V\in V.
\]
If $v\in H$, then $Jv\in H$, and since $dz_i|_H=0$,
\[
dz_i(Jv)=0=i\,dz_i(v).
\]
If $v\in V$, the standard fiberwise action of $J$ gives
\[
dz_i(Jv)=i\,dz_i(v).
\]
By linearity,
\[
dz_i(Jv)=i\,dz_i(v)
\]
for all $v\in TU$. Thus
\[
\bar\partial_J z_i=0,
\]
so $z_i$ is $J$-holomorphic.
\end{proof}

\begin{remark}
The lemma uses only the two special properties of an $\mathcal R$-compatible $J$ stated in Definition~\ref{def:AK}. These are the properties isolated in \cite{TMZ18}: the normal directions are modeled on the standard complex structure fiberwise, and the projection to each stratum is $J$-holomorphic. This is also consistent with the general discussion in \cite{MO}: for an arbitrary almost complex structure, finding local $J$-holomorphic defining functions is overdetermined and is controlled by the Nijenhuis tensor. In the almost K\"ahler data above, the Nijenhuis tensor of $J$ has image in the tangent directions of the strata, so the normal direction obstructions vanish; see \cite{TMZ18}.
\end{remark}

\section{Conclusion}

\begin{proof}[Proof of Theorem~\ref{thm:main}]
By the regularization existence theorem, Theorem~\ref{thm:reg}, after passing along a path of cohomologous symplectic forms we may assume that $D$ admits an $\omega_1$-regularization $\mathcal R$ and that there exists
\[
(\omega_1,\mathcal R,J)\in AK(M,D).
\]
The almost complex structure $J$ is $\omega_1$-compatible, hence $\omega_1$-tame.

Fix $p\in M$. If $p\notin D$, choose a neighborhood disjoint from $D$; then the required condition is vacuous with $r=0$. If $p\in D$, Lemma~\ref{lem:local} gives the desired smooth complex coordinates. Therefore the theorem follows.
\end{proof}

Thus the answer to the original question is yes: after replacing $\omega$ by a cohomologous symplectic form along a path preserving $D$ as a symplectic simple normal crossings divisor, there exists a tame almost complex structure for which the divisor is locally defined by $J$-holomorphic functions.

\begin{thebibliography}{TMZ21}

\bibitem{TMZ18}
M. Farajzadeh Tehrani, M. McLean, A. Zinger,
\href{https://www.numdam.org/item/10.1016/j.crma.2018.02.009.pdf}{numdam.org},
\textit{Singularities and semistable degenerations for symplectic topology},
C. R. Acad. Sci. Paris, Ser. I \textbf{356} (2018), 420--432.

\bibitem{TMZ21}
M. Farajzadeh Tehrani, M. McLean, A. Zinger,
\href{https://ar5iv.labs.arxiv.org/html/2112.13125}{ar5iv.labs.arxiv.org},
\textit{Normal crossings singularities for symplectic topology: structures},
arXiv:2112.13125.

\bibitem{FZ17}
M. Farajzadeh Tehrani, A. Zinger,
\href{https://ar5iv.labs.arxiv.org/html/1710.00224}{ar5iv.labs.arxiv.org},
\textit{Pseudoholomorphic curves relative to a normal crossings symplectic divisor: compactification},
arXiv:1710.00224.

\bibitem{NSF}
\href{https://par.nsf.gov/servlets/purl/10430629}{par.nsf.gov},
NSF public access repository copy of \cite{TMZ18}.

\bibitem{MO}
\href{https://mathoverflow.net/questions/154302/local-holomorphic-equations-for-symplectic-divisors}{mathoverflow.net},
\textit{Local holomorphic equations for symplectic divisors}.

\end{thebibliography}

\end{document}
```